\subsection{Fourier Orthogonality}

% **Role:** convergence of functions and analytic approximation.
%
% Trigonometric series, Fourier coefficients, convergence modes, orthogonality, Hilbert-space preview.


Fourier structure studies trigonometric series through orthogonality, projection, kernels, and convergence, with inner-product geometry reappearing throughout.

\begin{definition} \label{def:fourier-core-objects} \hlt{Core Fourier Objects}
\begin{enumerate}[label=(\roman*)]
\setlength{\itemsep}{0pt}
\item A \hlt{trigonometric polynomial} is a finite sum $a_0+\sum\limits_{n=1}^{N}(a_n\cos nx+b_n\sin nx)$, equivalently $\sum\limits_{n=-N}^{N}\alpha_ne^{inx}$.
\item If $f$ is Riemann integrable on $[-\pi,\pi]$, its \hlt{Fourier coefficients} are $\widehat{f}(m)\defeq \frac{1}{2\pi}\int_{-\pi}^{\pi}f(x)e^{-imx}\,dx$, $m\in\Z$, and $f(x)\sim_{\mathrm{coef}}\sum\limits_{m=-\infty}^{\infty}\widehat{f}(m)e^{imx}$ records coefficients, not convergence.
\item A sequence $(\phi_n)$ on $[a,b]$ is \hlt{orthonormal} if $\int_a^b \phi_n(x)\overline{\phi_m(x)}\,dx=0\quad(n\neq m)$ and $\int_a^b\abs{\phi_n(x)}^2\,dx=1$. Its Fourier coefficients are $c_n\defeq\int_a^b f(t)\overline{\phi_n(t)}\,dt$.
\end{enumerate}
\hlt{Key Insight}: Fourier coefficients are orthonormal coordinates; convergence is a separate analytic question.
\end{definition}

\begin{definition} \label{def:period-one-fourier-normalization} \hlt{Period-One Fourier Normalization}\\
Let $C(\R/\Z;\C)$ denote the continuous functions $f:\R\to\C$ satisfying $f(x+1)=f(x)$ for all $x$. \\
For $f,g\in C(\R/\Z;\C)$, set $\innerprod[\mathbb{T}]{f,g}\defeq \int_0^1 f(x)\overline{g(x)}\,dx$; $\norm{f}_{L^2(\mathbb{T})}\defeq \sqrt{\innerprod[\mathbb{T}]{f,f}}$; and $e_n(x)\defeq e^{2\pi i n x}$ for $n\in\Z$.\\
The \hlt{period-one Fourier coefficient} is $\widehat{f}_{\mathbb{T}}(n)\defeq \innerprod[\mathbb{T}]{f,e_n}=\int_0^1 f(x)e^{-2\pi i n x}\,dx$.\\
A period-one trigonometric polynomial is a finite sum $\sum_{n=-N}^{N}c_ne_n$.
\end{definition}

\begin{perspective} \label{per:period-one-versus-two-pi-fourier} \hlt{Two Equivalent Fourier Clocks}\\
The period-one and $2\pi$-periodic conventions differ by $t=2\pi x$; only normalizing constants change.\\
\hlt{Key Insight}: Changing Fourier clocks changes constants, not the projection geometry.
\end{perspective}

\begin{proposition} \label{prop:period-one-characters-orthonormal} \hlt{Characters Are an Orthonormal System}\\
For $m,n\in\Z$, $\innerprod[\mathbb{T}]{e_n,e_m}=1$ if $n=m$ and $0$ otherwise.\\
Hence, if $P=\sum\limits_{n=-N}^{N}c_ne_n$, then $c_n=\innerprod[\mathbb{T}]{P,e_n}$ and $\norm{P}_{L^2(\mathbb{T})}^2=\sum_{n=-N}^{N}\abs{c_n}^2$.
\begin{proof}
If $n=m$, the integrand is $1$. If $n\neq m$, then $\int_0^1e^{2\pi i(n-m)x}\,dx=0$.\\
The coefficient and norm identities are the finite orthonormal-coordinate identities from Theorems \ref{thm:orthonormal-basis-coordinates} and \ref{thm:parseval-identity-finite-dimensional}, applied inside the finite-dimensional span of $e_{-N},\ldots,e_N$.\\
\hlt{Structural Commentary}: Fourier coefficients are orthonormal coordinates in a finite trigonometric subspace.\\
\hlt{Key Insight}: The exponential characters make cross terms vanish.
\end{proof}
\end{proposition}

